\chapter{职业有害因素的识别与评价}
\setcounter{chapter}{10}
\chapterintro{
\textbf{【教学大纲要求】}

\imp{掌握：}职业卫生调查（调查类别、调查步骤）；作业环境监测（有害作业分级评价、作业场所空气中化学物质的检测、空气中有害物质测定方法）；生物监测（概念、特点、指标、工作程序、正常参比值和生物接触限值）；健康监护（概念、健康检查、健康监护档案、健康状况分析）；职业流行病学调查；职业性有害因素的危险度评定（内容）；职业病与工伤致残鉴定（功能能力评价、致残程度鉴定）。

\imp{熟悉：}实验研究；生物材料检测；职业卫生调查示例；危险度评定示例。
}

\section{职业卫生调查}

\subsection{调查类型}

\begin{enumerate}[leftmargin=*]
    \item \imp{基本情况调查}：了解企业基本职业卫生状况
    \item \imp{专题调查}：针对特定职业危害因素的深入调查
    \item \imp{应急性调查}：发生急性职业中毒或安全事故时进行
\end{enumerate}

\subsection{调查步骤}

\begin{enumerate}[leftmargin=*]
    \item 准备阶段（查阅资料、制定方案、准备仪器）
    \item 现场调查阶段（了解生产工艺、识别有害因素、采样监测）
    \item 资料整理和分析阶段
    \item 撰写调查报告
\end{enumerate}

\section{作业环境监测}

\subsection{有害化学物质的检测}

\begin{enumerate}[leftmargin=*]
    \item \imp{采样方法}：个体采样（个体采样器）和区域采样（定点采样）
    \item \imp{测定方法}：比色法、紫外可见分光光度法、原子吸收光谱法、气相/液相色谱法
\end{enumerate}

\subsection{职业接触限值（OEL）}

\begin{defbox}
\textbf{最高容许浓度（MAC）}：工作地点空气中任何一次有代表性的采样测定均不得超过的浓度。适用于急性危害较大的物质。

\textbf{时间加权平均容许浓度（PC-TWA）}：以时间为权数规定的8h工作日的平均容许接触水平。

\textbf{短时间接触容许浓度（PC-STEL）}：在遵守PC-TWA前提下容许15min短时间接触的浓度。

\textbf{制定依据：}有害物质的理化特性资料、动物实验资料、人体毒理学资料、现场职业卫生调查资料、流行病学调查资料。
\end{defbox}

\section{生物监测}

\begin{defbox}
\textbf{生物监测}：在规定的时间间隔内，系统地采集人体生物材料（血、尿、呼出气等），测定其中化学物质或其代谢产物的含量，或测定非损害性生化效应指标，以评价人体的内暴露水平和健康危险。

\textbf{特点：}
\begin{enumerate}[leftmargin=*]
    \item 反映机体总的内暴露剂量（考虑所有吸收途径和个体差异）
    \item 反映已吸收的量（而非环境中存在的量）
    \item 可反映机体对有害物质的代谢转化能力
    \item 可与健康效应指标相结合
\end{enumerate}

\textbf{生物监测指标：}
\begin{itemize}[leftmargin=*]
    \item 接触指标：血铅、尿汞、尿酚等
    \item 效应指标：ZPP、$\delta$-ALA、胆碱酯酶活性等
\end{itemize}

\textbf{生物接触限值（BEI）}：工人接触有害物质时，生物材料中该物质或其代谢产物或效应指标的容许限值。
\end{defbox}

\section{职业流行病学调查}

在职业卫生中的应用：
\begin{enumerate}[leftmargin=*]
    \item 阐明职业性损害在人群中的分布、发生和发展规律
    \item 发现职业性有害因素对健康的影响
    \item 为制定、修订职业卫生标准和职业病诊断标准提供依据
    \item 评价职业卫生和职业病防治工作质量及其预防措施效果
\end{enumerate}

\section{职业有害因素的危险度评价}

\begin{keybox}
\textbf{危险度评价（Risk Assessment）}：对职业有害因素损害健康的可能性和程度进行系统、科学的评价。

\textbf{内容：}
\begin{enumerate}[leftmargin=*]
    \item \imp{危害鉴定}：判断某因素是否对健康有害（定性）
    \item \imp{暴露评价}：评估暴露的强度、频率和持续时间
    \item \imp{剂量-反应关系评价}：量化剂量与效应关系（核心）
    \item \imp{危险度特征分析}：综合前三阶段，判断危害可能性大小
\end{enumerate}

\textbf{危险度管理}：根据危险度评价结果，结合经济、技术、社会等因素，制定和执行控制措施的过程。
\end{keybox}

\section{职业健康监护}

\begin{defbox}
\textbf{职业健康监护}：以预防为目的，根据劳动者职业接触史，通过定期或不定期的医学健康检查和健康相关资料收集，连续监测劳动者健康状况，分析健康变化与职业病危害因素的关系，及时采取干预措施。

\textbf{内容：}(1)就业前健康检查（发现职业禁忌证）；(2)定期健康检查（早期发现职业损害）；(3)离岗时健康检查；(4)应急健康检查。
\end{defbox}

% ----- 复习题 -----
\section{复习题}

\begin{enumerate}[leftmargin=*, itemsep=8pt]
    \item \textbf{【名词解释】}职业接触限值；MAC；PC-TWA；PC-STEL
    \item \textbf{【名词解释】}生物监测；生物接触限值（BEI）；职业健康监护
    \item \textbf{【名词解释】}危险度评价；危险度管理
    \item \textbf{【简答题】}简述职业卫生调查的类型和步骤。
    \item \textbf{【简答题】}简述生物监测的特点和常用指标。
    \item \textbf{【简答题】}简述职业有害因素危险度评价的四个阶段。
    \item \textbf{【简答题】}简述职业健康监护的内容。
\end{enumerate}

\subsection*{参考答案要点}

\begin{enumerate}[leftmargin=*]
    \item \textbf{职业接触限值}：为保护劳动者健康而规定的职业接触有害因素不应超过的限量值。\textbf{MAC}：任何一次采样均不得超过的最高浓度。\textbf{PC-TWA}：8h工作日时间加权平均容许浓度。\textbf{PC-STEL}：15min短时间接触容许浓度。
    \item \textbf{生物监测}：系统采集人体生物材料测定化学物或其代谢产物含量以评价内暴露水平。\textbf{BEI}：生物材料中容许的限值。\textbf{职业健康监护}：通过医学检查连续监测劳动者健康状况。
    \item \textbf{危险度评价}：对有害因素损害健康的可能性和程度的系统评价。\textbf{危险度管理}：基于评价结果制定和执行控制措施。
    \item 三类：基本情况调查、专题调查、应急性调查。步骤：准备→现场调查→资料整理分析→撰写报告。
    \item 特点：反映总内暴露剂量（所有途径）、反映已吸收量、反映个体代谢差异、可与健康效应指标结合。指标：接触指标（血铅、尿汞）、效应指标（ZPP、ChE活性）。
    \item 四个阶段：危害鉴定（定性判断是否对健康有害）→暴露评价（评估暴露强度、频率和持续时间）→剂量-反应关系评价（核心，量化关系）→危险度特征分析（综合判断危害可能性）。
    \item 就业前健康检查（发现禁忌证）、定期健康检查（早期发现损害）、离岗时健康检查、应急健康检查。
\end{enumerate}

\newpage

% =====================================================================
%  CHAPTER 11: 职业有害因素的监督与控制
% =====================================================================
\newpage